% Setting up the document with beamer class for slides
\documentclass{beamer}
\usetheme{Copenhagen}
\usecolortheme{seahorse}

% Including necessary packages for UTF-8 and Vietnamese support
\usepackage[utf8]{inputenc}
\usepackage[vietnamese]{babel}
\usepackage{lmodern}
\usepackage{booktabs}
\usepackage{multirow}
\usepackage{graphicx} % Cần cho \resizebox
\usepackage{xcolor}
\usepackage{amsmath} % Cần cho \text{}

% Defining title, author, and date
\title{Báo cáo Hiệu suất ASR (5 Tập Dữ liệu)}
\author{}
\date{\today}

% Starting the document
\begin{document}

% Title slide
\begin{frame}
    \titlepage
\end{frame}

% Slide 1: Bối cảnh
\begin{frame}{Bối cảnh}
    \begin{itemize}
        \item \textbf{Chunkformer (đóng, đang dùng):} 
            \begin{itemize}
                \item WER $23-27\%$, tốt nhưng \textbf{không thể tinh chỉnh}.
            \end{itemize}
        \item \textbf{Fast-Conformer (mở, huấn luyện sẵn VietSpeech):}
            \begin{itemize}
                \item WER $35-39\%$, kém.
            \end{itemize}
        \item \textbf{Hướng đi:} Tinh chỉnh 2 bước:
            \begin{itemize}
                \item v1: Nhãn giả (pseudo label).
                \item v2: Nhãn giả $\to$ Dữ liệu huấn luyện VPB.
            \end{itemize}
    \end{itemize}
\end{frame}

% Slide 2: Kết quả WER (5 tập dữ liệu)
\begin{frame}{Kết quả WER (5 tập dữ liệu)}
    \centering
    \resizebox{\textwidth}{!}{ % Co bảng vừa với slide
        \begin{tabular}{lccccp{2cm}}
            \toprule
            \textbf{Tập dữ liệu} & \textbf{Chunkformer} & \textbf{v0 (vietspeech)} & \textbf{v1 (pseudo)} & \textbf{v2 (pseudo$\to$VPB)} & \textbf{Ghi chú} \\
            \midrule
            standard\_test\_2 & 0.2499 & 0.3546 & 0.3040 & 0.2420 & \text{$ 70\%$ overlap} \\
            standard\_test & 0.1613 & 0.3378 & 0.2582 & \textbf{0.4294} & \text{Nhỏ (29 mẫu)} \\
            next\_day\_test\_debug & 0.2076 & 0.3414 & 0.2687 & 0.2649 & \text{Test độc lập} \\
            vpb\_right2\_train & 0.2420 & 0.3633 & 0.2956 & 0.2456 & \text{Tập huấn luyện} \\
            vpb\_right2\_valid & 0.2696 & 0.3902 & 0.3240 & 0.2821 & \text{Tập xác thực} \\
            \bottomrule
        \end{tabular}
    }
\end{frame}

% Slide 3: Phân tích
\begin{frame}{Phân tích}
    \begin{itemize}
        \item \textbf{Huấn luyện sẵn}: Kém xa Chunkformer (WER $ 0.35$--$0.39$).
        \item \textbf{Tinh chỉnh v1}: Cải thiện mạnh (WER giảm $ 0.07$--$0.10$).
        \item \textbf{Tinh chỉnh v2}: Tiệm cận Chunkformer ở \textbf{tập chính (validation, next\_day)}.
        \item \textbf{Lưu ý}:
            \begin{itemize}
                \item WER thấp ở \textbf{standard\_test\_2} do \textbf{rò rỉ dữ liệu train} (\text{70\% overlap}).
                \item WER thấp ở \textbf{train} thể hiện fine-tune hội tụ, \textbf{không đại diện cho kết quả với new user}.
            \end{itemize}
    \end{itemize}
\end{frame}

% Slide 4: Insight chính
\begin{frame}{Insight chính}
    \begin{itemize}
        \item \textbf{Vấn đề cốt lõi}: Giọng người dùng đa dạng \& nhiều nhiễu:
            \begin{itemize}
                \item Giọng vùng miền, tốc độ, nói lắp, từ địa phương.
                \item Nhiễu môi trường, tín hiệu cuộc gọi kém.
            \end{itemize}
        \item \textbf{Thách thức}: Tạo mô hình \textbf{robust for unseen data (user voice)}.
        \item \textbf{Giải pháp}: Gán nhãn thủ công $\to$ tập dữ liệu chuẩn nội bộ, liên tục cập nhật voice của new user 
        \item \textbf{Tương lai}: Thu thập user voice ngay từ khi làm hợp đồng vay nợ $\to$ xác thực user, đòi nợ, nhắc nợ tự động qua voice-call
    \end{itemize}
\end{frame}

% Slide 5: Giá trị gán nhãn
\begin{frame}{Giá trị gán nhãn}
    \begin{itemize}
        \item \textbf{Điều kiện cần} để tinh chỉnh và nâng cao hiệu năng.
        \item \textbf{Tài sản dữ liệu dài hạn}:
            \begin{itemize}
                \item Phân tích cuộc gọi, xác thực user, đánh giá chất lượng dịch vụ.
                \item Huấn luyện mô hình NLP.
                \item Phát triển Callbot (tương tác tự động qua voice).
            \end{itemize}
        \item \textbf{Đầu tư một lần}, tái sử dụng nhiều vòng.
    \end{itemize}
\end{frame}

% Slide 6: Tính cấp bách
\begin{frame}{Tính cấp bách}
    \begin{itemize}
        \item \textbf{Tín dụng AWS}: 15,000 USD (hết hạn 20/09/2025).
        \item \textbf{Hạn gán nhãn}: 15/09/2025.
        \item \textbf{Nhu cầu}: 100 nhân viên, 500h làm thêm ($ 100$ triệu VND).
    \end{itemize}
\end{frame}

% Slide 7: Kết luận
\begin{frame}{Kết luận}
    \begin{itemize}
        \item \textbf{Không gán nhãn}: Mô hình dừng ở 70--75\%.
        \item \textbf{Có gán nhãn}: Đạt \textbf{75--80\%+}, sở hữu \textbf{dữ liệu chuẩn chiến lược}.
        \item \textbf{Đầu tư làm thêm}: Hợp lý để tận dụng tín dụng AWS \& phát triển Callbot bền vững.
    \end{itemize}
\end{frame}

% Ending the document
\end{document}